% ======================================================================
%  CAPÍTULO 19 — GLOSARIO (parte final de la parte V)
% ======================================================================

\chapter{Glosario}\label{ch:glosario}

\begin{roadmap}
Este glosario reúne, en orden alfabético, los términos que el libro
usa con un significado técnico preciso. Cada entrada recuerda en una
o dos líneas la definición y remite al capítulo donde se desarrolla.
Hace además las veces de índice: quien busque «¿dónde se habló de
X?» puede venir aquí primero.
\end{roadmap}

\medskip
\noindent Los nombres de algoritmos y los parámetros aparecen con su
notación del texto; los capítulos se citan por su número de
referencia.

\begin{description}[leftmargin=2.6em,style=nextline,itemsep=2.5pt,parsep=0pt]

  \item[\textbf{Acortamiento} (\emph{shortening})]
        Técnica que fija a cero $s$ coordenadas de un código
        $(n,k)$; en ISD transporta la dureza entre instancias con
        pérdida $\le0{,}8$ bits, y es la cara amable de tomar
        $n<q$ (caps.~\ref{ch:isd} y~\ref{ch:estructura}).

  \item[\textbf{AES-128 / AES-256}]
        Cifrado simétrico de referencia. Grover reduce su búsqueda
        de clave a $2^{64}$ y $2^{128}$: el «listón de 128 bits»
        que traducen las categorías (caps.~\ref{ch:pqc}
        y~\ref{ch:isd}).

  \item[\textbf{Alfa} ($\alpha$; orden del soporte)]
        Lista ordenada $\alpha_0,\dots,\alpha_{q-1}$ de todos los
        elementos de $\GF q$; el soporte del código es
        $L=\{\alpha_0,\dots,\alpha_{n-1}\}$. Su \emph{orden} es
        secreto y se sortea con $\textsf{FieldOrdering}$
        (cap.~\ref{ch:keygen}).

  \item[\textbf{Anulación de orden superior} (\emph{higher-order
        vanishing})]
        Familia de distinguidores de 2025--2026 que evalúa el
        código contra funciones de correlación de orden creciente;
        de ella salen el \emph{holdout} y las estimaciones de
        GIJS/Weis (caps.~\ref{ch:estructura} y~\ref{ch:2026}).

  \item[\textbf{Benes, red de}]
        Representación de la permutación $\pi$ (el orden $\alpha$)
        con $2m-1$ etapas de $2^{m-1}$ conmutadores:
        $\lceil(2m-1)2^{m-4}\rceil$ bytes (5\,888 para $m=12$);
        es como la clave secreta guarda $\alpha$
        (cap.~\ref{ch:keygen}).

  \item[\textbf{Berson, ataque de reacción de}]
        Con un descifrado que avisa del fallo, $n$ consultas
        $x\oplus e_j$ revelan el soporte del error. El rechazo
        implícito del KEM lo desactiva (cap.~\ref{ch:1978}).

  \item[\textbf{BSI} (\emph{Bundesamt für Sicherheit in der
        Informationstechnik})]
        Oficina federal alemana de seguridad. Su nota del
        1-oct-2026 desaconseja Classic McEliece en nuevos
        despliegues tras las estimaciones de 2026 y recomienda
        ML-KEM, FrodoKEM o HQC (cap.~\ref{ch:2026}).

  \item[\textbf{Canal lateral}]
        Ataque a la implementación (tiempos, cachés, accesos)
        que no rompe el esquema abstracto; el ataque de tiempos de
        Shoup contra la decodificación es el ejemplo clásico
        (caps.~\ref{ch:kem} y~\ref{ch:2026}).

  \item[\textbf{Categorías (NIST)}]
        Escala que traduce «romper el esquema» a costes tipo AES:
        categoría~1 $\approx2^{143}$ operaciones de puerta
        (AES-128), categoría~5 $\approx2^{272}$ (AES-256). Classic
        McEliece se certifica con el modelo de costes que carga la
        memoria (caps.~\ref{ch:parametros} y~\ref{ch:isd}).

  \item[\textbf{Clase lateral}]
        Conjunto $y+C$ de vectores con el mismo síndrome que $y$;
        decodificar es elegir en la clase el representante de peso
        mínimo (cap.~\ref{ch:lineales}).

  \item[\textbf{Clave equivalente}]
        Objeto (por ejemplo $g$ y el orden $\alpha$, o un
        decodificador del código público) con el que se descifra
        todo sin ser la clave secreta literal; su obtención es una
        rotura total (caps.~\ref{ch:keygen}, \ref{ch:seguridad}
        y~\ref{ch:estructura}).

  \item[\textbf{Código de Goppa} ($\Gamma_2(L,g)$)]
        Código $\{c:\sum_j c_j/(x+\alpha_j)\equiv0\pmod g\}$, con
        $g$ de grado $t$ sin raíces en $L$; la versión binaria
        (subcuerpo) con $g$ irreducible tiene distancia
        $\ge2t+1$: es el código que usa Classic McEliece
        (cap.~\ref{ch:goppa}).

  \item[\textbf{Código lineal}]
        Subespacio de $\Ftwo^n$ descrito por una matriz generadora
        $G$ o una de paridad $\Hmat$; su distancia mínima $d$ es
        el peso del codeword no nulo más ligero
        (cap.~\ref{ch:lineales}).

  \item[\textbf{Conjetura de dureza (SD en media)}]
        «Para todo algoritmo subexponencial, el éxito sobre
        instancias de la distribución del esquema es
        despreciable»: el supuesto en el que descansa la seguridad
        (caps.~\ref{ch:seguridad} y~\ref{ch:isd}).

  \item[\textbf{Cuerpo finito} ($\GF q$ con $q=2^m$; $\Ftwo$)]
        El terreno donde viven $g$, $\alpha$ y los síndromes;
        $\Ftwo$ es el cuerpo base y $\GF q$ el de extensión
        (cap.~\ref{ch:cuerpos}).

  \item[\textbf{Delta} ($\delta$; semilla)]
        Los $32$ bytes que siembran toda la generación: de ella se
        derivan deterministamente el código, el orden y el secreto
        de rechazo $s$; la clave secreta completa es regenerable
        desde $\delta$ (cap.~\ref{ch:keygen}).

  \item[\textbf{Decode / Encode}]
        Primitivas del KEM: $\textsf{Encode}(e)=\hat{\Hmat}e$ y
        $\textsf{Decode}(C)$, que corrige hasta $t$ errores y solo
        acepta vectores de peso exactamente $t$
        (caps.~\ref{ch:encapdecap} y~\ref{ch:labkem}).

  \item[\textbf{Despliegues}]
        Implementaciones reales y usos del esquema: SUPERCOP y
        \texttt{cmimpl}, liboqs y PQClean, y los despliegues
        híbridos (Mullvad, Rosenpass, PQConnect), todos anteriores
        al aviso del BSI de 2026 (cap.~\ref{ch:2026}).

  \item[\textbf{Distancia de Hamming; distancia mínima} ($d$)]
        El peso de la diferencia entre dos vectores; $d$ es el
        mínimo peso no nulo de un codeword. El teorema de Goppa da
        $d\ge2t+1$ y de ahí la decodificación única
        (caps.~\ref{ch:lineales} y~\ref{ch:goppa}).

  \item[\textbf{Distinguidor}]
        Algoritmo que decide «¿Goppa de verdad o matriz aleatoria?»
        con ventaja no despreciable. Erosiona el supuesto de
        indistinguibilidad de la clave pública, pero no descifra
        por sí solo (caps.~\ref{ch:seguridad}--\ref{ch:2026}).

  \item[\textbf{Dumer, algoritmo de}]
        Mejora de Prange (1991) que precomputa columnas y
        amplifica la probabilidad de éxito por iteración; primer
        escalón relevante tras Prange (cap.~\ref{ch:isd}).

  \item[\textbf{$e$ (vector de error)}]
        Vector de peso exactamente $t$ que cifra:
        $C=\hat{\Hmat}e$. Su soporte es el «mensaje-error» del
        esquema (caps.~\ref{ch:1978}--\ref{ch:encapdecap}).

  \item[\textbf{Escalera de ataques}]
        La sucesión que baja la cotización de ISD para cada
        conjunto: Prange (1962) $\to$ Lee--Brickell/Leon (1988)
        $\to$ Stern/Dumer (1989/91) $\to$ MMT/BJMM (2011/12)
        $\to$ May--Mezerov (2015) $\to$ May--Ozerov (2018)
        $\to$ Esser--Bellini y Carrier (2022) $\to$ las revisiones
        de 2026 (cap.~\ref{ch:isd}).

  \item[\textbf{Espacio de criptogramas} ($P$)]
        La imagen de los errores de peso $t$ bajo $\textsf{Encode}$
        ($\binom{n}{t}$ elementos; $|P|>2^{456}$ en
        \texttt{348864}); aparece en la cota ROM
        (cap.~\ref{ch:kem}).

  \item[\textbf{Faugère, ataque de}]
        Recuperación algebraica de la clave por sistemas
        polinómicos con interrupción temprana (XLS de Faugère et
        al.); hoy limitado a $t\le20$: holgura enorme frente a los
        $t$ reales (caps.~\ref{ch:estructura} y~\ref{ch:2026}).

  \item[\textbf{FieldOrdering}]
        Algoritmo de la especificación que convierte $32q$ bits en
        el orden uniforme $\alpha$; el choque de enteros provoca
        reintento determinista (cap.~\ref{ch:keygen};
        proposición~\ref{prop:field-ordering}).

  \item[\textbf{FIPS 203/204/205}]
        Los estándares post-cuánticos del NIST (ML-KEM, ML-DSA,
        SLH-DSA); el ecosistema del que Classic McEliece es
        alternativa y complemento (caps.~\ref{ch:pqc}
        y~\ref{ch:2026}).

  \item[\textbf{FixedWeight}]
        Algoritmo que muestrea errores uniformes de peso exacto
        $t$ con trozos de $16$ bits y rechazo; su uniformidad es
        una proposición (cap.~\ref{ch:encapdecap};
        proposición~\ref{prop:fixedweight-uniform}).

  \item[\textbf{FO (Fujisaki--Okamoto)}]
        Transformación genérica de PKE a KEM en el ROM; Classic
        McEliece usa una versión simplificada (sin reencriptación,
        con rechazo implícito y secretos más pequeños)
        (cap.~\ref{ch:kem}).

  \item[\textbf{FrodoKEM}]
        KEM de retículas (LWE) considerado conservador; una de las
        alternativas que el BSI recomienda mientras se revisa la
        TR-02102-1 (cap.~\ref{ch:2026}).

  \item[\textbf{$G'$ (matriz pública de McEliece)}]
        La generadora disfrazada $G'=SGP$; su versión sistemática
        se reduce a la clave pública $T$ (caps.~\ref{ch:1978}
        y~\ref{ch:niederreiter}).

  \item[\textbf{$g$ (polinomio de Goppa)}]
        Polinomio mónico de grado $t$ sobre $\GF q$ sin raíces en
        $L$; el secreto central del código, sorteado con
        $\textsf{Irreducible}$ (caps.~\ref{ch:goppa}
        y~\ref{ch:keygen}).

  \item[\textbf{Gaussiana (eliminación)}]
        Procedimiento para llevar $\hat{\Hmat}$ a la forma
        $(\mathrm{I}_{mt}\,|\,T)$; también el corazón de las
        iteraciones de ISD (caps.~\ref{ch:labgoppa}
        y~\ref{ch:isd}).

  \item[\textbf{Grover, algoritmo de}]
        Búsqueda cuántica con aceleración cuadrática en el
        exponente: AES-256 $\to2^{128}$; en ISD, la «mitad» del
        exponente, con la latencia como límite realista
        (caps.~\ref{ch:pqc} y~\ref{ch:parametros}).

  \item[\textbf{Hamming, código de} {($[7,4,3]$)}]
        El ejemplo recurrente para síndromes y corrección: matriz
        de paridad con las siete columnas no nulas y decodificación
        de un error (cap.~\ref{ch:lineales}).

  \item[\textbf{Harvest now, decrypt later}]
        Grabar hoy el tráfico cifrado para descifrarlo cuando
        exista el ordenador cuántico; la razón de migrar antes de
        que sea urgente (cap.~\ref{ch:pqc}).

  \item[\textbf{Híbrido}]
        Combinación de un KEM post-cuántico con uno clásico (por
        ejemplo X25519); los despliegues reales de Classic
        McEliece son híbridos, lo que conserva la mitad clásica
        ante roturas (caps.~\ref{ch:pqc} y~\ref{ch:2026}).

  \item[\textbf{Holgura} (\emph{tightness})]
        Pérdida entre una cota de seguridad y el juego real: el
        QROM pierde la raíz cuadrada del éxito OW-CPA
        (cap.~\ref{ch:kem}).

  \item[\textbf{HQC}]
        KEM de códigos cuasi-cíclicos seleccionado por el NIST en
        2025; la otra alternativa que el BSI pone por delante de
        McEliece (caps.~\ref{ch:parametros} y~\ref{ch:2026}).

  \item[\textbf{Irreducible}]
        Algoritmo de la especificación que convierte $16t$ bits en
        un $g$ mónico irreducible uniforme; su demostración es la
        proposición de uniformidad (cap.~\ref{ch:keygen}).

  \item[\textbf{ISD (decodificación por conjuntos de
        información)}]
        Familia de algoritmos de decodificación genérica (Prange,
        Stern, MMT, BJMM, May--Ozerov, \dots) que fija la
        seguridad clásica de los conjuntos (cap.~\ref{ch:isd}).

  \item[\textbf{ISO/IEC 18033-2 (Amd.2)}]
        Norma con los $4$ conjuntos $\times$ $4$ variantes $=16$
        perfiles de Classic McEliece (jun-2026); descarta
        \texttt{348864} (caps.~\ref{ch:parametros} y~\ref{ch:2026}).

  \item[\textbf{Juguete (el)}]
        La instancia diminuta usada en los laboratorios:
        $m=4$ ($q=16$), $n=16$, $t=2$, elaborada paso a paso en el
        capítulo~\ref{ch:labgoppa} y completa en
        el~\ref{ch:labkem}; \texttt{toy.py}
        (apéndice~\ref{app:codigo}) reproduce sus números.

  \item[\textbf{$K$ (clave del KEM)}]
        Los $32$ bytes que encapsula y desencapsula el KEM:
        $H(1,e,C)$ en el éxito y $H(0,s,C)$ en el rechazo
        (caps.~\ref{ch:kem} y~\ref{ch:labkem}).

  \item[\textbf{KAT (vectores de prueba)}]
        Archivos de la especificación con claves, criptogramas y
        secretos de referencia; comparables byte a byte porque todo
        el proceso (incluidos los reintentos) es determinista
        (cap.~\ref{ch:keygen}).

  \item[\textbf{KEM (encapsulamiento de claves)}]
        Primitiva \textsf{Encap}/\textsf{Decap} que produce una
        clave compartida a partir de una clave pública; la que
        Classic McEliece deriva de su PKE (cap.~\ref{ch:kem}).

  \item[\textbf{Latencia (cuántica)}]
        Límite práctico de Grover: las iteraciones se serializan y
        la raíz cuadrada sobreestima al atacante real; de ahí que
        la «mitad» del exponente se compare con cautela
        (caps.~\ref{ch:parametros} y~\ref{ch:2026}).

  \item[\textbf{Lee--Brickell}]
        Primer algoritmo de ISD con lista (1988): fija $p$
        posiciones «con error» y busca en ellas; semilla de la
        línea Stern (cap.~\ref{ch:isd}).

  \item[\textbf{Leon}]
        Variante de lista por información canónica (1988),
        hermana de Lee--Brickell (cap.~\ref{ch:isd}).

  \item[\textbf{Localizador (polinomio $\sigma$)}]
        Polinomio cuyas raíces señalan las posiciones del error:
        $\sigma(x)=\prod_{j\in\supp(e)}(x+\alpha_j)$; Patterson lo
        recupera del síndrome (caps.~\ref{ch:goppa}
        y~\ref{ch:labgoppa}).

  \item[\textbf{$m$}]
        Exponente del cuerpo: $q=2^m$. Classic McEliece usa
        $m=12$ ($q=4096$) o $m=13$ ($q=8192$); la seguridad
        estructural crece $\approx8$ bits por unidad de $m$
        (caps.~\ref{ch:parametros} y~\ref{ch:estructura}).

  \item[\textbf{MatGen}]
        Subalgoritmo de generación de claves: forma
        semi-sistemática de las $mt$ primeras columnas, con
        reintentos deterministas ($\delta\leftarrow\delta'$)
        (cap.~\ref{ch:keygen}).

  \item[\textbf{May--Mezerov / May--Ozerov}]
        Ataques que llevan MMT/BJMM y la búsqueda de vecinos a
        ISD; May--Ozerov (2018) cotiza \texttt{348864} en
        $2^{140{,}8}$ (memoria libre) (cap.~\ref{ch:isd}).

  \item[\textbf{McEliece (criptosistema)}]
        El cifrado de 1978: $x=uG'+z$ con la generadora disfrazada
        y error de peso $t$; quien conoce el código lo decodifica
        con Patterson (cap.~\ref{ch:1978}).

  \item[\textbf{Memoria (coste de la)}]
        En ISD, las tablas grandes dominan el coste real: cargarla
        añade decenas de bits (hasta $33{,}5$ en \texttt{6960119})
        y decide las categorías (caps.~\ref{ch:isd}
        y~\ref{ch:2026}).

  \item[\textbf{ML-KEM}]
        El KEM post-cuántico estandarizado por el NIST (FIPS 203);
        referencia de comparación para la criptografía de códigos
        (caps.~\ref{ch:pqc} y~\ref{ch:2026}).

  \item[\textbf{Modelo de costes}]
        La convención con la que se cuenta el trabajo de un ataque
        (puertas, bits, operaciones vectoriales, memoria,
        latencia); las columnas «libre» y «con cargo» de las
        tablas (caps.~\ref{ch:isd}, \ref{ch:estructura}
        y~\ref{ch:2026}).

  \item[\textbf{NCP (decodificación / palabra más cercana)}]
        El problema hermano de SD: dado $y$, hallar el codeword de
        peso mínimo a distancia de $y$; el descifrado de McEliece
        es una instancia (caps.~\ref{ch:lineales}
        y~\ref{ch:seguridad}).

  \item[\textbf{Niederreiter (reformulación)}]
        Versión dual basada en síndromes: $C=e(H')^{\top}$;
        equivalente a McEliece por la reducción de la
        proposición~\ref{prop:nied-reduccion}
        (cap.~\ref{ch:niederreiter}).

  \item[\textbf{NIST}]
        Organizador de la competición post-cuántica (2016--2024);
        estandarizó ML-KEM/ML-DSA/SLH-DSA, aplazó su decisión
        sobre McEliece y seleccionó HQC como KEM extra
        (caps.~\ref{ch:pqc} y~\ref{ch:2026}).

  \item[\textbf{OW-CPA}]
        Juego de seguridad del PKE: recuperar $e$ dado $C$; la
        raíz que el QROM divide por dos en su cota
        (caps.~\ref{ch:kem} y~\ref{ch:seguridad}).

  \item[\textbf{Parámetros (conjuntos de)}]
        Los cinco del estándar: \texttt{348864}, \texttt{460896},
        \texttt{6688128}, \texttt{6960119}, \texttt{8192128}
        ($m=12$ o $13$, $t=64$--$128$); sus tablas de tamaños y
        seguridad (caps.~\ref{ch:parametros}--\ref{ch:2026}).

  \item[\textbf{Paridad (matriz de, $\Hmat$)}]
        Matriz cuyas filas ortogonalizan el código
        ($GH^{\top}=0$); en forma sistemática
        $\hat{\Hmat}=(\mathrm{I}_{mt}\,|\,T)$
        (caps.~\ref{ch:lineales} y~\ref{ch:keygen}).

  \item[\textbf{Patterson (algoritmo de)}]
        Decodificador de los códigos de Goppa: forma el síndrome
        $S$, resuelve la ecuación clave (Euclides/Berlekamp) y
        halla las raíces del localizador; $O(t^2)$ operaciones en
        $\GF q$ más inversiones (caps.~\ref{ch:goppa},
        \ref{ch:labgoppa} y~\ref{ch:encapdecap}).

  \item[\textbf{pcf / pc / f (variantes)}]
        Sufijos de los conjuntos: \texttt{f} = generación sin
        reintentos visibles (forma $(32,64)$); \texttt{pc} =
        confirmación de texto claro ($C=(C_0,C_1)$ con
        $C_1=H(2,e)$); \texttt{pcf} = ambas (cap.~\ref{ch:parametros}).

  \item[\textbf{Permutación ($P$, $\pi$)}]
        El disfraz de coordenadas, $(\vv P)_j=\vv_{\sigma(j)}$; su
        orden es el secreto que FieldOrdering sortea
        (caps.~\ref{ch:labgoppa} y~\ref{ch:keygen}).

  \item[\textbf{Peso} ($\wt$)]
        Número de posiciones no nulas. Los parámetros del esquema
        son pesos: $\wt(e)=t$ y $\wt(z)\le t$
        (caps.~\ref{ch:lineales} y~\ref{ch:1978}).

  \item[\textbf{Pi} ($\pi$)]
        La permutación que resulta de ordenar los enteros de
        FieldOrdering; se guarda comprimida en la red de Benes
        (cap.~\ref{ch:keygen}).

  \item[\textbf{PKE (cifrado de clave pública)}]
        El cifrado subyacente (McEliece/Niederreiter) del que el
        KEM se deriva mediante la transformación del
        capítulo~\ref{ch:kem}.

  \item[\textbf{PQC (criptografía post-cuántica)}]
        El campo entero: algoritmos diseñados para resistir
        atacantes con ordenador cuántico
        (caps.~\ref{ch:pqc} y~\ref{ch:2026}).

  \item[\textbf{Prange}]
        ISD original (1962): elegir $k$ columnas y comprobar si el
        error cae fuera de ellas; iteración barata y probabilidad
        pequeña. La referencia contra la que se mide todo lo demás
        (cap.~\ref{ch:isd}).

  \item[\textbf{QROM}]
        Modelo del oráculo aleatorio cuántico: la cota de FO
        pierde la raíz cuadrada del éxito OW-CPA
        (caps.~\ref{ch:kem} y~\ref{ch:estructura}).

  \item[\textbf{Recuperación de clave} (\emph{key recovery})]
        El objetivo más fuerte de la taxonomía estructural:
        obtener $g$, el orden $\alpha$ o un decodificador; a él
        apuntan las cifras de 2026 ($2^{140}$--$2^{148}$,
        $2^{99}$--$2^{102}$ según contabilidad)
        (caps.~\ref{ch:estructura} y~\ref{ch:2026}).

  \item[\textbf{Rechazo implícito}]
        \textsf{Decap} nunca devuelve $\bot$: entrega una clave
        pseudoaleatoria $H(0,s,C)$; elimina el oráculo
        observacional de los protocolos (caps.~\ref{ch:kem}
        y~\ref{ch:encapdecap}).

  \item[\textbf{ROM}]
        Modelo del oráculo aleatorio clásico; la cota aditiva de
        la transformación de Fujisaki--Okamoto (cap.~\ref{ch:kem}).

  \item[\textbf{$s$ (secreto de rechazo)}]
        Cadena uniforme de $n$ bits dentro de la clave secreta;
        la clave de la rama de rechazo es $H(0,s,C)$
        (caps.~\ref{ch:keygen} y~\ref{ch:encapdecap}).

  \item[\textbf{SD (Syndrome Decoding)}]
        El problema: dado $(H,y)$ con $y=He$, hallar $e$ de peso
        $\le t$. NP-duro en el peor caso; duro en media (conjetura)
        sobre la distribución del esquema (caps.~\ref{ch:seguridad}
        y~\ref{ch:isd}).

  \item[\textbf{Semi-sistemática (forma, $(\mu,\nu)$)}]
        Truco de \textsf{MatGen}: permutar hasta que las $mt$
        primeras columnas sean independientes;
        $(\mu,\nu)=(32,64)$ en las variantes \texttt{f}
        (cap.~\ref{ch:keygen}).

  \item[\textbf{SHAKE256}]
        El XOF de la especificación: expande semillas, genera
        $g$ y $\alpha$, y deriva las claves $K$ y de rechazo
        (caps.~\ref{ch:keygen} y~\ref{ch:labkem}).

  \item[\textbf{Shor, algoritmo de}]
        Rompe RSA y ECC en tiempo polinómico: la razón de ser de
        toda la migración post-cuántica (cap.~\ref{ch:pqc}).

  \item[\textbf{Shoup, ataque de tiempos de}]
        Ataque de canal lateral clásico contra implementaciones de
        decodificación; la motivación histórica de decodificar en
        tiempo constante (cap.~\ref{ch:kem}).

  \item[\textbf{Síndrome}]
        $S(y)=\Hmat y^{\top}$; en el código de Goppa también
        $S(x)=\sum_j c_j/(x+\alpha_j)$ con columnas
        $\sigma_j=(x+\alpha_j+1)/g(\alpha_j)$; su espacio tiene
        dimensión $n-k$ (caps.~\ref{ch:lineales}
        y~\ref{ch:labgoppa}).

  \item[\textbf{Soporte} ($\supp$)]
        El conjunto de posiciones no nulas de un vector. El error
        del KEM es uniforme sobre los $\binom{n}{t}$ soportes de
        tamaño $t$ (caps.~\ref{ch:lineales}
        y~\ref{ch:encapdecap}).

  \item[\textbf{Stern (algoritmo de)}]
        ISD con cruce de listas (1989): el primer gran salto de la
        escalera ($22$ bits sobre Prange en \texttt{348864})
        (cap.~\ref{ch:isd}).

  \item[\textbf{$t$}]
        Grado del polinomio de Goppa y peso exacto del error
        ($t=64$--$128$ en los conjuntos); fija la distancia
        $\ge2t+1$ y el trabajo de Patterson
        (caps.~\ref{ch:goppa} y~\ref{ch:parametros}).

  \item[\textbf{$\Tmat$ (clave pública)}]
        La mitad derecha de $\hat{\Hmat}=(\mathrm{I}_{mt}|T)$: la
        clave pública, de $mt\lceil k/8\rceil$ bytes
        ($261\,120$ a $1\,357\,824$) (caps.~\ref{ch:keygen}
        y~\ref{ch:parametros}).

  \item[\textbf{Terremoto (2025--2026)}]
        El periodo de distinguidores, estimaciones y contabilidades
        en disputa que movió la conversación pública sin romper el
        esquema (caps.~\ref{ch:estructura} y~\ref{ch:2026}).

  \item[\textbf{TII}]
        Retos públicos de instancias Goppa pequeñas ($n=214$)
        usados para probar heurísticas estructurales; no existe
        (a octubre de 2026) un cómputo de núcleos contra un
        conjunto real (cap.~\ref{ch:2026}).

  \item[\textbf{Traza} ($\Tr$)]
        El funcional $\Tr(\beta)=\sum_{i=0}^{m-1}\beta^{2^i}$;
        criterios de irreducibilidad, raíces cuadradas y
        linealidad sobre $\Ftwo$ (cap.~\ref{ch:cuerpos}).

  \item[\textbf{Vectores (notación)}]
        Las cadenas binarias se imprimen con la posición $0$ a la
        izquierda; al serializar se empaqueta en little-endian
        (posición $j$ en el bit $j\bmod8$ del byte $j/8$)
        (caps.~\ref{ch:labkem} y~\ref{app:codigo}).

  \item[\textbf{Weis (ataques de 2026)}]
        Recuperación heurística de clave vía \emph{holdout},
        con cifras entre $2^{89}$ y $2^{128}$ según la
        contabilidad, y la regla «$\approx8$ bits por unidad de
        $m$» (caps.~\ref{ch:estructura} y~\ref{ch:2026}).

\end{description}

\medskip
\noindent\textit{Fin del glosario. Los apéndices que siguen cierran el
libro: el código comentado del juguete (apéndice~\ref{app:codigo}),
la cronología de medio siglo de historia y las tablas del cuerpo
$\FF4$.}
